\chapter*{အမှာစာ}
\addcontentsline{toc}{chapter}{အမှာစာ}

ဒီစာအုပ်က neural network တွေအကြောင်းကို သင်္ချာအခြေခံ ဦးစားပေး မိတ်ဆက်ပေးမဲ့ စာအုပ်တစ်အုပ်အနေနဲ့
ရည်ရွယ်ပြီး ရေးသားထားတာ ဖြစ်ပါတယ်။
neural network တွေအကြောင်း deep learning အခြေခံတွေကိုတော့ ဒီစာအုပ် တစ်အုပ်ထဲ
ဖတ်ပြီး အကုန်ရသွားဖို့ မဖြစ်နိုင်ပါဘူး။ ဒီစာအုပ်သာမက အခြားလိုအပ်တဲ့ စာအုပ်တွေ၊ ဗီဒီယိုတွေ
အပြင် အွန်လိုင်း သင်တန်းတွေပါ ပူးတွဲလေ့လာမှ အခြေခံကို ပြည့်ပြည့်စုံစုံ မြင်မှာဖြစ်ပါတယ်။

\bigskip

ကျွန်တော်ရဲ့ ရည်ရွယ်ချက်က သင်္ချာအခြေခံကောင်းတွေကို စောစောစီးစီး သိသွားအောင်ပါပဲ။
တစ်ချို့သော သင်္ချာသဘောတရားတွေဟာ မိတ်ဆွေအတွက် အနည်းငယ် နားလည်ရခက်နေမယ်။
ဒါပေမဲ့ စိတ်မပျက်ပါနဲ့။ ဒီလိုတွေရှိတယ်။ ဒီလို သင်္ချာတွေက တကယ်အရေးကြီးပါလားဆိုတာကို
သိစေချင်လို့ ထည့်ပေးထားတာပါ။
စောစောနားလည်သွားရင် သင်္ချာကို ပိုပြီး အလေးအနက်ထား လေ့လာဖြစ်သွားအောင် ရည်ရွယ်ပါတယ်။

\bigskip

အကြံပေးရမယ်ဆိုရင် ဒီစာအုပ်ကို ဖြည့်စွက်စာအုပ်တစ်အုပ်အဖြစ် ဖတ်စေချင်ပါတယ်။
ဒီစာအုပ်က မိတ်ဆွေအတွက် neural network အကြောင်း ပထမဦးဆုံးပြောပြမဲ့ စာအုပ် မဖြစ်သင့်ပါဘူး။
Youtube ကနေဖြစ်ဖြစ်၊ ကျောင်း(သို့)သင်တန်းကနေ neural network အကြောင်း အနည်းအကျဥ်းသိထားမှာသာ ပိုပြီးလက်ခံနိုင်မှာဖြစ်ပါတယ်။

\bigskip

နောက်ပြီး ဒီစာအုပ်ကို ဖတ်ပြီး စိတ်မဝင်စားတာပဲ ဖြစ်ဖြစ်၊ Deep learning နဲ့ neural network တွေကို အရမ်းခက်တယ်လို့ ထင်သွားရင် မိတ််ဆွေမှားတာမဟုတ်ပါဘူး။
ကျွန်တော်ရဲ့ ရေးသားပုံ၊ ရှင်းပြပုံ အားနည်းချက်ကြောင့်သာ ဖြစ်မှာပါ။
ဒီစာအုပ်ထက် ပိုပြီး သင့်တော်မဲ့ လေ့လာစရာတွေကို ရှာကြည့်ပါ။
ယခုခေတ်က AI ခေတ်ဖြစ်သွားပြီမို့ AI အကြောင်း၊ AI ပညာရပ်ကို သိနိုင်သလောက်သိအောင်
ကြိုးစားဖို့ အကြံပေးချင်ပါတယ်။

\vspace{2em}
\noindent\textit{Light}
